\chapter{Introduction}
\label{ch:intro}


\section{研究背景}

隨著行動通訊網路邁向第六代（6G），網路架構的複雜度與服務多樣性正面臨挑戰。5G標準引入了網路切片（Network Slicing）的概念，允許在同一套實體基礎設施上承載特性迥異的邏輯網路，例如以高吞吐量為導向的增強行動寬頻（Enhanced Mobile Broadband, eMBB）與以超低延遲為核心的超可靠低延遲通訊（Ultra-Reliable and Low Latency Communications, URLLC）。然而，如何在多租戶、多切片環境中實現快速、精確且自動化的端到端配置，仍是電信產業與學術界共同面臨的核心難題。

傳統的網路管理方式高度依賴人工介入與靜態配置，維運人員需要深入理解底層設備的廠商專屬介面，方能完成核網(Core)與無線接取網(RAN)的協調配置。這種方式不僅耗時，更難以因應 6G 所要求的動態服務等級協議（Service Level Agreement, SLA）調整。以具體場景為例，當一組工業無人機的控制訊號需要從 eMBB 切片臨時升級至 URLLC 切片，以滿足更嚴苛的封包延遲預算（Packet Delay Budget, PDB）時，傳統流程往往耗時數分鐘乃至數十分鐘，遠超實際應用所能容忍的上限。

為回應上述挑戰，3GPP TS 28.312 規範提出了意圖驅動管理（Intent-driven Management）框架，允許上層應用程式（rApp）以高階、聲明式的語意描述網路期望狀態，而無需關心底層實現細節。然而，此類標準框架的落地實踐，尤其是在涵蓋 RAN 與 Core 的完整端到端閉環控制上，目前仍缺乏足夠具體的開源參考實作 \cite{3gpp_ts28312}。

與此同時，雲原生技術正與電信領域深度融合。Nephio 專案以 Kubernetes 資源模型（Kubernetes Resource Model, KRM）為基礎，透過聲明式 GitOps 工作流程管理網路功能的部署與配置，提供了一套有別於傳統運維支援系統（OSS）的新型電信自動化平台。Nephio 的 Porch 套件管理系統允許對設備配置進行版本控制、審查與差異化部署，從而將電信網路的配置管理與軟體開發的持續整合實踐相結合 \cite{nephio}。

\section{研究動機}

行動網路的管理體系正經歷從命令式（Imperative）向宣告式（Declarative）的根本性轉變。傳統模式在處理跨 Core 與 RAN 的協調時，往往因缺乏跨域語意理解而導致配置流程遲滯。隨著 6G 網路切片技術對 SLA 動態調整的要求日益嚴苛，傳統管理架構已無法因應即時的切片遷移需求。雖然 3GPP TS 28.312 定義了意圖驅動管理框架，試圖透過「期望狀態」簡化管理複雜度，但該規範多停留在理論層次，在跨域閉環控制與實際部署流程上仍缺乏具體的實作參照。

現有的雲原生編排工具與電信開源專案之間存在技術斷層。以 Nephio 為例，其架構重心在於套件的生命週期管理，本身並不具備將高階業務意圖（如延遲、頻寬需求）自動轉譯為底層領域參數，如 PRB Ratio（Physical Resource Block Ratio）或 5QI（5G QoS Identifier）的語意轉譯引擎。此功能缺失導致 rApp 下發的意圖無法直接作用於底層網路元件。此外，如 srsRAN 等主流 RAN 模擬軟體，目前仍缺乏與 Nephio 深度整合的 Operator，且常用的 UE（User Equipment）模擬器（如 srsue）多以獨立軟體形式存在，缺乏標準化部署方案。這種工具鏈的碎片化限制了全自動化實驗環境的構建，也使得跨域意圖編排的驗證流程難以完整串接。

基於上述在語意轉譯引擎、跨域自動化工作流以及雲原生模擬環境整合上的落差，建立一個能夠串接異質南向介面並具備宣告式意圖處理能力的架構，已成為推進電信自動化的關鍵路徑。本研究之動機在於解決現有框架在動態切片配置上的遲滯問題，藉由填補概念性規範與實體配置管道之間的斷層，探索如何透過雲原生技術實現高度解耦且可驗證的端到端意圖驅動編排。

\section{研究目標與主要貢獻}

針對前述研究動機，本研究之首要目標為開發一套遵循 3GPP TS 28.312 規範的雲原生端到端意圖編排系統。本研究透過 Kubernetes Operator 模式實作宣告式 R1 介面，藉由設計 E2E Intent Orchestrator 核心組件，將抽象的 SLA 需求轉化為可執行的領域配置指令。此系統之建構旨在達成從意圖提交、參數轉譯到無線電資源生效的完整自動化管道。

為支撐上述架構，本研究在技術實作與系統整合方面達成以下貢獻。首先，本研究開發了專為 Nephio 體系設計的 srsran-operator，使 srsRAN 網路功能得以整合至雲原生編排流程，並達成與 free5GC 的自動化連接。其次，本研究實作了 srsran-helm 部署方案，將 srsue 封裝為符合雲原生規範的應用組件。此工作將原本獨立運行的模擬工具轉化為可受 Kubernetes 調度的標準化資源，為端到端意圖驅動系統提供了一個可自動化執行且易於擴展的測試平台。

本研究的另一項主要貢獻在於實作了跨網域的混合南向介面。編排器透過整合 Nephio Porch 的宣告式工作流，實現對 RAN 參數（如 PRB Ratio 與優先權）的非同步更新；同時，系統透過調用 free5GC WebConsole 的應用程式介面（API），同步完成 Core 的使用者註冊與 QoS 參數配置。這種混合編排架構證明了在維持各網域異質性的前提下，透過單一意圖入口達成跨域閉環控制的可行性。

最後，本研究透過實驗驗證了所提系統在資源隔離與部署效率上的有效性。實驗數據顯示，在實體資源受限的環境下，系統能根據意圖指令精確分配 PRB，確保多種切片共存時的服務品質穩定。同時，經由系統性量測，從使用者提交意圖至無線電資源配置實際生效的端到端延遲成功控制在 30 至 36 秒內。這些具體成果展示了意圖驅動管理在提升電信維運自動化方面的實踐價值。
